% see the original template for more detail about bibliography, tables, etc: https://www.overleaf.com/latex/templates/handout-design-inspired-by-edward-tufte/dtsbhhkvghzz

\documentclass{tufte-handout}

%\geometry{showframe}% for debugging purposes -- displays the margins

\usepackage{amsmath}

\usepackage{hyperref}

\usepackage{fancyhdr}

\usepackage{hanging}

\hypersetup{colorlinks=true,allcolors=[RGB]{97,15,11}}

\fancyfoot[L]{\emph{History of Media Studies}, vol. 6, 2026}


% Set up the images/graphics package
\usepackage{graphicx}
\setkeys{Gin}{width=\linewidth,totalheight=\textheight,keepaspectratio}
\graphicspath{{graphics/}}

\title[Critical Communication]{Critical Communication: A Memoir} % longtitle shouldn't be necessary

% The following package makes prettier tables.  We're all about the bling!
\usepackage{booktabs}

% The units package provides nice, non-stacked fractions and better spacing
% for units.
\usepackage{units}

% The fancyvrb package lets us customize the formatting of verbatim
% environments.  We use a slightly smaller font.
\usepackage{fancyvrb}
\fvset{fontsize=\normalsize}

% Small sections of multiple columns
\usepackage{multicol}

% Provides paragraphs of dummy text
\usepackage{lipsum}

% These commands are used to pretty-print LaTeX commands
\newcommand{\doccmd}[1]{\texttt{\textbackslash#1}}% command name -- adds backslash automatically
\newcommand{\docopt}[1]{\ensuremath{\langle}\textrm{\textit{#1}}\ensuremath{\rangle}}% optional command argument
\newcommand{\docarg}[1]{\textrm{\textit{#1}}}% (required) command argument
\newenvironment{docspec}{\begin{quote}\noindent}{\end{quote}}% command specification environment
\newcommand{\docenv}[1]{\textsf{#1}}% environment name
\newcommand{\docpkg}[1]{\texttt{#1}}% package name
\newcommand{\doccls}[1]{\texttt{#1}}% document class name
\newcommand{\docclsopt}[1]{\texttt{#1}}% document class option name


\begin{document}

\begin{titlepage}

\begin{fullwidth}
\noindent\LARGE\emph{Book review
} \hspace{88mm}\includegraphics[height=1cm]{logo3.png}\\
\noindent\hrulefill\\
\vspace*{1em}
\noindent{\Huge{\emph{Critical Communication: A Memoir}\par}}

\vspace*{1.5em}

\noindent\LARGE{John A. Lent}\par\marginnote{\emph{Critical Communication: A Memoir}, reviewed by John A. Lent, \emph{History of Media Studies} 6 (2026), \href{https://doi.org/https://doi.org/10.32376/d895a0ea.7efe94dc}{https://doi.org/10.32376/d895a0ea.7efe94dc}.} \vspace*{0.75em}
\vspace*{0.5em}
\noindent{{\large\emph{Temple University}, \href{mailto:john.lent@gmail.com}{john.lent@gmail.com}\par}} 

\end{fullwidth}

\vspace*{1em}

\noindent Vincent Mosco\marginnote{\href{https://creativecommons.org/licenses/by-nc/4.0/}{\includegraphics[height=0.5cm]{by-nc.png}}}. \emph{Critical Communication: A Memoir}. 266 pp. University of Westminster Press, 2024. \$22.99 (paper).
\vspace*{2em}



\newthought{Full disclosure:}  I wrote one of the testimonials for this book, and its
author was a personal acquaintance. Over a span of forty-five years, I
knew Vinny Mosco as a colleague at Temple University in the early stages
of his career; as one of the co-organizers of the Union for Democratic
Communications in the early 1980s; as the supervisor of his wife, Cathy
McKercher's, master of journalism degree; as a contributor to
\emph{Global Productions}, which I edited with Gerald Sussman; as a
fellow author in books on political economy and/or critical
communication; and as one of twelve key thinkers in critical
communications scholarship that I selected to be included in a 2015 book
that I co-edited with Michelle A. Amazeen. With that out of the way, we
can talk about this book.

From the beginning, in the preface, Mosco alerts us that his memoir is
different from others: shorter, to avoid overwriting and bloating with
details; generous, in that, as he writes, ``it is not my intention to
settle scores, pick fights, or dress up my life with the lurid details
that might briefly enliven the text only to diminish the more core of
the story'' (xxi); and structured, around the theme of critical
communication, with occasional dips into major times in his life.


Mosco's accounts of growing up in an Italian family in a three-room
tenement apartment in New York's Little Italy, and how it shaped his
lifelong values and attitudes, hit a chord with this reviewer, also of
Italian descent, who grew up amidst small town poverty. Mosco's
recounting of ``living without,'' inventing and organizing his own
games, taking ``sponge baths'' in the kitchen sink for lack of a tub,
and enduring cold winter nights with just the heat 



\enlargethispage{2\baselineskip}

\vspace*{2em}

\noindent{\emph{History of Media Studies}, vol. 6, 2026}


\end{titlepage}


\noindent generated by the
kitchen stove, is presented with some levity, while also showing how
these hardships early on carved out his future research and activism in
the areas of poverty, social class, and labor.

The reportage is honest and upfront, as was Mosco, presented to provide
background and amplify key moments, decisions, and individuals (good
guys and scoundrels) that had an impact on his personal and professional
life. He goes to some length to explain a couple uncles' and aunts'
connections to the Mafia, and his father's abhorrence of the Cosa
Nostra, and his fight with the draft board to obtain conscientious
objector status, as well as his acceptances to a prestigious secondary
school, Georgetown University, and Harvard University. With
characteristic humbleness, Mosco talks about his friendships at
Georgetown with future president Bill Clinton and roommate Joe Mitchell,
nephew of later notorious attorney general John Mitchell.

As the heading of the second chapter suggests, Mosco's experiences at
Georgetown and Harvard gave him opportunities first to ``learn the
canon'' and then to challenge it. A course, ``The Problem of God,'' led
him to a final decision to reject formal religion, and sharp divisions
in Georgetown's School of Foreign Service, in which he was enrolled,
made him realize that he was not cut out to be a diplomat or State
Department employee, which led him to transfer to the College of Arts
and Sciences, and pursue a life of teaching and research. By this time,
he had learned ``to navigate a variety of intellectual landscapes'' by
using the dimensions of ``historical (where does the issue come from?),
epistemological (how best to study it?), ontological (what level of
reality does it occupy---material, symbolic?).'' Mosco was to use this
model in his redevelopment of the theory of political economy
communication three decades later.

Throughout, his work is meticulously arranged by formulae of his own
making, based on lessons he discovered studying and/or working (and
disagreeing) with stellar educators and researchers such as Carroll
Quigley, Daniel Bell, and Anthony Oettinger, and through his eagerness
to move into ever-new research territory and his voracious reading of
classical works in a number of disciplines. Along with his exploration
of critical communication and his first-hand account of the history of
political economy, Mosco raises several fresh and the unconventional
topics, like the cloud, ``smart cities,'' utopia, myth, and the digital
sublime.

Mosco, and later also McKercher, lived most of the time by unorthodox
precepts. Among these were that ``the fun of life took precedence over
even our best thought-out plans,'' that it is important to learn
patience and live in the present, and that tenure, which Mosco
relinquished at three universities, is for people who put job security
first and foremost. Mosco understood the appeal of security, but in his
case, when he had to choose between a career and life {[}experience{]},
he chose the latter.

As with all people, Mosco's life was marked with ups and downs. He
recounts times of depression, which he fought through therapy, thankful
enough to dedicate his \emph{To The Cloud} to his therapist, Louise
LaPlante. He also documents his fears of losing both his wife and one of
his daughters to breast cancer and the sadness of losing his brother,
Anthony, to drug addiction. Such hardships were countered by celebratory
moments. He and McKercher won the Association for Education in
Journalism and Mass Communication Professional Freedom and
Responsibility Award, joining previous honorees Noam Chomsky, I. F.
Stone, Studs Terkel, among others; \emph{The Digital Sublime} was
awarded the 2005 Olson Prize for outstanding book in rhetorical and
cultural theory; and a Manhattan street was named Mosco, to commemorate
the vast public service given by Vinny's father, Frank, and the success
attained by Mosco's former students, particularly the five PhD
candidates he supervised while at Queen's University. In fact, about
seven pages of the last chapter of this memoir relate to the successful
careers of his students.

Mosco also takes time to compare the academic scenes of his birth
country and his adopted one, pointing out that Canadian universities are
much more supportive with grants and more tolerant of political
disagreement. He hones in on his stay at Temple University, where the
then-dean of communications feared and refused to support critical
communications, rejecting Dallas Smythe's request for an extension of
his contract, accusing Mosco of turning his colleague Janet Wasko into a
radical, and demeaning the critical communication thinkers by saying
that when he hired Mosco, he thought ``{[}he{]} was only getting one of
{[}them{]}.'' I was a Temple professor at the time and remember how we
were labeled a ``Marxist cell,'' which more than likely was the dean's
view.

\emph{Critical Communication: A Memoir} is a riveting read, full of
philosophies of life and work, funny and instructive anecdotes, a
rundown of the evolution of the fields of critical communications and
the political economy of communications, and tidbits about the trappings
of academia, good and bad. The book is recommended for those who have
experienced poverty, class bias, and racism, to refresh their memories,
and for the more privileged to learn from; for colleagues, students,
friends, and fellow scholars worldwide, to recall Mosco's admirable
traits; for upper-level undergraduate and graduate classes, to learn
about many facets of communications theory and practice; and for other
activists, to make them aware that as long as the life and visions of
Vinny Mosco are remembered, they are not alone in their struggle.

\enlargethispage{2\baselineskip}

\end{document}